% ---------------------------------------------------------------------------
% A posteriori error estimation via self-convergence.
% Intended location: % ---------------------------------------------------------------------------
% A posteriori error estimation via self-convergence.
% Intended location: % ---------------------------------------------------------------------------
% A posteriori error estimation via self-convergence.
% Intended location: % ---------------------------------------------------------------------------
% A posteriori error estimation via self-convergence.
% Intended location: \input{self_convergence_section} after the Accuracy
% Study (\cref{seq:accuracy}).
% ---------------------------------------------------------------------------
\subsection{A Posteriori Error Estimation via Self-Convergence}
\label{seq:self_convergence}

The accuracy studies above rely on a sequential reference solution, which is
unavailable in production---computing it would negate the parallel speedup.
\method\ can instead estimate its own truncation error. Let $u^{(0)}$ denote
the truncated parallel solution of \cref{eq:truncated_form}. By
\cref{eq:global_error_sum}, its error is exactly the sum of the omitted
thermal tails $v_{i\to j}$, and these tails can be recovered incrementally:
iteration $k$ extends the correction window of every retained dependency by
$\Delta$ segments, computes the omitted tails on the newly exposed window
with the source-off operator (\cref{eq:correction_recursion}), and superposes
them to form the refined iterate $u^{(k)}$. No source-on solve is repeated
and no DAG is rebuilt. Writing $d_k = u^{(k)} - u^{(k-1)}$, each omitted
tail enters exactly one iteration, so \cref{eq:global_error_sum} telescopes
into
\begin{equation}
u^{\mathrm{exact}} - u^{(0)} \;=\; \sum_{k\ge 1} d_k ,
\label{eq:telescoping}
\end{equation}
up to contributions of predecessors never connected in the DAG, which are
below the calibrated threshold by \cref{eq:edge_criterion} and negligible in
all our experiments.

Two consequences follow. First,
$\|u^{(k)} - u^{(0)}\|$ converges to the true error of $u^{(0)}$ as the
iterations exhaust the omitted tails---an error estimate that requires no
reference solution---while the decay of $\|d_k\|$ signals convergence: the
iteration stops once $\|d_k\|$ falls below the accuracy target. The
increments decay because successive windows capture tails at increasingly
large elapsed times, which the contractive source-off evolution
(\cref{eq:stability_contraction}) monotonically attenuates; the decay is
monotone in practice provided the window increment $\Delta$ spans the local
revisit scale of the scan path. Second, the
procedure repairs what it measures: each $u^{(k)}$ is itself a more accurate
solution, so a run that detects a tolerance violation has already computed
its remedy. All norms below are the relative $L^2$ metric of
\cref{seq:accuracy}, maximized over the sampled snapshots.

We validate the estimator on the Bull path with the $10^{-4}$ configuration
of \cref{tab:acc_1e4}, partitioned across 32 ranks, where the many partition
boundaries drive the observed error slightly above the target.
\cref{tab:self_convergence} reports six iterations with $\Delta = 4$: each
$\|d_k\|$ reproduces the true error of the previous iterate to three or more
significant digits, the converged estimate
$\|u^{(6)} - u^{(0)}\| = 1.0767\times10^{-4}$ equals the true production
error to all reported digits, one iteration restores the run to its target,
and six reduce the error by three orders of magnitude.

\begin{table}[h]
\centering
\caption{Self-convergence iterations on the Bull path ($10^{-4}$
configuration of \cref{tab:acc_1e4}, 32 ranks, $\Delta=4$). True errors are
measured against a sequential reference used for validation only; $\|d_k\|$
is computed by the parallel solver alone.}
\label{tab:self_convergence}
\begin{tabular}{ccc}
\hline
Iterate & True max relative $L^2$ error & $\|d_k\|$ \\
\hline
$u^{(0)}$ & $1.0767\times10^{-4}$ & --- \\
$u^{(1)}$ & $3.6498\times10^{-5}$ & $1.0767\times10^{-4}$ \\
$u^{(2)}$ & $1.6113\times10^{-5}$ & $3.6497\times10^{-5}$ \\
$u^{(3)}$ & $9.9789\times10^{-6}$ & $1.6113\times10^{-5}$ \\
$u^{(4)}$ & $3.3293\times10^{-6}$ & $9.9778\times10^{-6}$ \\
$u^{(5)}$ & $8.9734\times10^{-7}$ & $3.3284\times10^{-6}$ \\
$u^{(6)}$ & $1.8131\times10^{-7}$ & $8.9678\times10^{-7}$ \\
\hline
\end{tabular}
\end{table}
 after the Accuracy
% Study (\cref{seq:accuracy}).
% ---------------------------------------------------------------------------
\subsection{A Posteriori Error Estimation via Self-Convergence}
\label{seq:self_convergence}

The accuracy studies above rely on a sequential reference solution, which is
unavailable in production---computing it would negate the parallel speedup.
\method\ can instead estimate its own truncation error. Let $u^{(0)}$ denote
the truncated parallel solution of \cref{eq:truncated_form}. By
\cref{eq:global_error_sum}, its error is exactly the sum of the omitted
thermal tails $v_{i\to j}$, and these tails can be recovered incrementally:
iteration $k$ extends the correction window of every retained dependency by
$\Delta$ segments, computes the omitted tails on the newly exposed window
with the source-off operator (\cref{eq:correction_recursion}), and superposes
them to form the refined iterate $u^{(k)}$. No source-on solve is repeated
and no DAG is rebuilt. Writing $d_k = u^{(k)} - u^{(k-1)}$, each omitted
tail enters exactly one iteration, so \cref{eq:global_error_sum} telescopes
into
\begin{equation}
u^{\mathrm{exact}} - u^{(0)} \;=\; \sum_{k\ge 1} d_k ,
\label{eq:telescoping}
\end{equation}
up to contributions of predecessors never connected in the DAG, which are
below the calibrated threshold by \cref{eq:edge_criterion} and negligible in
all our experiments.

Two consequences follow. First,
$\|u^{(k)} - u^{(0)}\|$ converges to the true error of $u^{(0)}$ as the
iterations exhaust the omitted tails---an error estimate that requires no
reference solution---while the decay of $\|d_k\|$ signals convergence: the
iteration stops once $\|d_k\|$ falls below the accuracy target. The
increments decay because successive windows capture tails at increasingly
large elapsed times, which the contractive source-off evolution
(\cref{eq:stability_contraction}) monotonically attenuates; the decay is
monotone in practice provided the window increment $\Delta$ spans the local
revisit scale of the scan path. Second, the
procedure repairs what it measures: each $u^{(k)}$ is itself a more accurate
solution, so a run that detects a tolerance violation has already computed
its remedy. All norms below are the relative $L^2$ metric of
\cref{seq:accuracy}, maximized over the sampled snapshots.

We validate the estimator on the Bull path with the $10^{-4}$ configuration
of \cref{tab:acc_1e4}, partitioned across 32 ranks, where the many partition
boundaries drive the observed error slightly above the target.
\cref{tab:self_convergence} reports six iterations with $\Delta = 4$: each
$\|d_k\|$ reproduces the true error of the previous iterate to three or more
significant digits, the converged estimate
$\|u^{(6)} - u^{(0)}\| = 1.0767\times10^{-4}$ equals the true production
error to all reported digits, one iteration restores the run to its target,
and six reduce the error by three orders of magnitude.

\begin{table}[h]
\centering
\caption{Self-convergence iterations on the Bull path ($10^{-4}$
configuration of \cref{tab:acc_1e4}, 32 ranks, $\Delta=4$). True errors are
measured against a sequential reference used for validation only; $\|d_k\|$
is computed by the parallel solver alone.}
\label{tab:self_convergence}
\begin{tabular}{ccc}
\hline
Iterate & True max relative $L^2$ error & $\|d_k\|$ \\
\hline
$u^{(0)}$ & $1.0767\times10^{-4}$ & --- \\
$u^{(1)}$ & $3.6498\times10^{-5}$ & $1.0767\times10^{-4}$ \\
$u^{(2)}$ & $1.6113\times10^{-5}$ & $3.6497\times10^{-5}$ \\
$u^{(3)}$ & $9.9789\times10^{-6}$ & $1.6113\times10^{-5}$ \\
$u^{(4)}$ & $3.3293\times10^{-6}$ & $9.9778\times10^{-6}$ \\
$u^{(5)}$ & $8.9734\times10^{-7}$ & $3.3284\times10^{-6}$ \\
$u^{(6)}$ & $1.8131\times10^{-7}$ & $8.9678\times10^{-7}$ \\
\hline
\end{tabular}
\end{table}
 after the Accuracy
% Study (\cref{seq:accuracy}).
% ---------------------------------------------------------------------------
\subsection{A Posteriori Error Estimation via Self-Convergence}
\label{seq:self_convergence}

The accuracy studies above rely on a sequential reference solution, which is
unavailable in production---computing it would negate the parallel speedup.
\method\ can instead estimate its own truncation error. Let $u^{(0)}$ denote
the truncated parallel solution of \cref{eq:truncated_form}. By
\cref{eq:global_error_sum}, its error is exactly the sum of the omitted
thermal tails $v_{i\to j}$, and these tails can be recovered incrementally:
iteration $k$ extends the correction window of every retained dependency by
$\Delta$ segments, computes the omitted tails on the newly exposed window
with the source-off operator (\cref{eq:correction_recursion}), and superposes
them to form the refined iterate $u^{(k)}$. No source-on solve is repeated
and no DAG is rebuilt. Writing $d_k = u^{(k)} - u^{(k-1)}$, each omitted
tail enters exactly one iteration, so \cref{eq:global_error_sum} telescopes
into
\begin{equation}
u^{\mathrm{exact}} - u^{(0)} \;=\; \sum_{k\ge 1} d_k ,
\label{eq:telescoping}
\end{equation}
up to contributions of predecessors never connected in the DAG, which are
below the calibrated threshold by \cref{eq:edge_criterion} and negligible in
all our experiments.

Two consequences follow. First,
$\|u^{(k)} - u^{(0)}\|$ converges to the true error of $u^{(0)}$ as the
iterations exhaust the omitted tails---an error estimate that requires no
reference solution---while the decay of $\|d_k\|$ signals convergence: the
iteration stops once $\|d_k\|$ falls below the accuracy target. The
increments decay because successive windows capture tails at increasingly
large elapsed times, which the contractive source-off evolution
(\cref{eq:stability_contraction}) monotonically attenuates; the decay is
monotone in practice provided the window increment $\Delta$ spans the local
revisit scale of the scan path. Second, the
procedure repairs what it measures: each $u^{(k)}$ is itself a more accurate
solution, so a run that detects a tolerance violation has already computed
its remedy. All norms below are the relative $L^2$ metric of
\cref{seq:accuracy}, maximized over the sampled snapshots.

We validate the estimator on the Bull path with the $10^{-4}$ configuration
of \cref{tab:acc_1e4}, partitioned across 32 ranks, where the many partition
boundaries drive the observed error slightly above the target.
\cref{tab:self_convergence} reports six iterations with $\Delta = 4$: each
$\|d_k\|$ reproduces the true error of the previous iterate to three or more
significant digits, the converged estimate
$\|u^{(6)} - u^{(0)}\| = 1.0767\times10^{-4}$ equals the true production
error to all reported digits, one iteration restores the run to its target,
and six reduce the error by three orders of magnitude.

\begin{table}[h]
\centering
\caption{Self-convergence iterations on the Bull path ($10^{-4}$
configuration of \cref{tab:acc_1e4}, 32 ranks, $\Delta=4$). True errors are
measured against a sequential reference used for validation only; $\|d_k\|$
is computed by the parallel solver alone.}
\label{tab:self_convergence}
\begin{tabular}{ccc}
\hline
Iterate & True max relative $L^2$ error & $\|d_k\|$ \\
\hline
$u^{(0)}$ & $1.0767\times10^{-4}$ & --- \\
$u^{(1)}$ & $3.6498\times10^{-5}$ & $1.0767\times10^{-4}$ \\
$u^{(2)}$ & $1.6113\times10^{-5}$ & $3.6497\times10^{-5}$ \\
$u^{(3)}$ & $9.9789\times10^{-6}$ & $1.6113\times10^{-5}$ \\
$u^{(4)}$ & $3.3293\times10^{-6}$ & $9.9778\times10^{-6}$ \\
$u^{(5)}$ & $8.9734\times10^{-7}$ & $3.3284\times10^{-6}$ \\
$u^{(6)}$ & $1.8131\times10^{-7}$ & $8.9678\times10^{-7}$ \\
\hline
\end{tabular}
\end{table}
 after the Accuracy
% Study (\cref{seq:accuracy}).
% ---------------------------------------------------------------------------
\subsection{A Posteriori Error Estimation via Self-Convergence}
\label{seq:self_convergence}

The accuracy studies above rely on a sequential reference solution, which is
unavailable in production---computing it would negate the parallel speedup.
\method\ can instead estimate its own truncation error. Let $u^{(0)}$ denote
the truncated parallel solution of \cref{eq:truncated_form}. By
\cref{eq:global_error_sum}, its error is exactly the sum of the omitted
thermal tails $v_{i\to j}$, and these tails can be recovered incrementally:
iteration $k$ extends the correction window of every retained dependency by
$\Delta$ segments, computes the omitted tails on the newly exposed window
with the source-off operator (\cref{eq:correction_recursion}), and superposes
them to form the refined iterate $u^{(k)}$. No source-on solve is repeated
and no DAG is rebuilt. Writing $d_k = u^{(k)} - u^{(k-1)}$, each omitted
tail enters exactly one iteration, so \cref{eq:global_error_sum} telescopes
into
\begin{equation}
u^{\mathrm{exact}} - u^{(0)} \;=\; \sum_{k\ge 1} d_k ,
\label{eq:telescoping}
\end{equation}
up to contributions of predecessors never connected in the DAG, which are
below the calibrated threshold by \cref{eq:edge_criterion} and negligible in
all our experiments.

Two consequences follow. First,
$\|u^{(k)} - u^{(0)}\|$ converges to the true error of $u^{(0)}$ as the
iterations exhaust the omitted tails---an error estimate that requires no
reference solution---while the decay of $\|d_k\|$ signals convergence: the
iteration stops once $\|d_k\|$ falls below the accuracy target. The
increments decay because successive windows capture tails at increasingly
large elapsed times, which the contractive source-off evolution
(\cref{eq:stability_contraction}) monotonically attenuates; the decay is
monotone in practice provided the window increment $\Delta$ spans the local
revisit scale of the scan path. Second, the
procedure repairs what it measures: each $u^{(k)}$ is itself a more accurate
solution, so a run that detects a tolerance violation has already computed
its remedy. All norms below are the relative $L^2$ metric of
\cref{seq:accuracy}, maximized over the sampled snapshots.

We validate the estimator on the Bull path with the $10^{-4}$ configuration
of \cref{tab:acc_1e4}, partitioned across 32 ranks, where the many partition
boundaries drive the observed error slightly above the target.
\cref{tab:self_convergence} reports six iterations with $\Delta = 4$: each
$\|d_k\|$ reproduces the true error of the previous iterate to three or more
significant digits, the converged estimate
$\|u^{(6)} - u^{(0)}\| = 1.0767\times10^{-4}$ equals the true production
error to all reported digits, one iteration restores the run to its target,
and six reduce the error by three orders of magnitude.

\begin{table}[h]
\centering
\caption{Self-convergence iterations on the Bull path ($10^{-4}$
configuration of \cref{tab:acc_1e4}, 32 ranks, $\Delta=4$). True errors are
measured against a sequential reference used for validation only; $\|d_k\|$
is computed by the parallel solver alone.}
\label{tab:self_convergence}
\begin{tabular}{ccc}
\hline
Iterate & True max relative $L^2$ error & $\|d_k\|$ \\
\hline
$u^{(0)}$ & $1.0767\times10^{-4}$ & --- \\
$u^{(1)}$ & $3.6498\times10^{-5}$ & $1.0767\times10^{-4}$ \\
$u^{(2)}$ & $1.6113\times10^{-5}$ & $3.6497\times10^{-5}$ \\
$u^{(3)}$ & $9.9789\times10^{-6}$ & $1.6113\times10^{-5}$ \\
$u^{(4)}$ & $3.3293\times10^{-6}$ & $9.9778\times10^{-6}$ \\
$u^{(5)}$ & $8.9734\times10^{-7}$ & $3.3284\times10^{-6}$ \\
$u^{(6)}$ & $1.8131\times10^{-7}$ & $8.9678\times10^{-7}$ \\
\hline
\end{tabular}
\end{table}
